\scanpage{52}
\mspage{6}

à un élément. Alors pour tte flèche $f:x\longrightarrow y$ dans $X$, \deleted{$\lambda\unsure{\mapsto|-|\to} f$}
\[
\lambda\longmapsto f\lambda:\quad H_x\overset{\sim}{\longrightarrow} H_y
\]
est bijectif, (i.e les $H_x$ forment un syst. local sur $X$).

$4^\circ)$ Supposons, \deleted{\unsure{que}} en plus des conditions de a) b) c) et celles de $3^\circ)$, que $X$ soit $1$-connexe. Alors le système local des $H_x$ est \inserted{donc} constant, \deleted{\unsure{donc|dessus}} en particulier l'opération \inserted{à droite} de $M=\operatorname{Hom}(e,e)$ sur $H_x$ est constante, \deleted{i.e $\lambda\in$} et par suite l'
\[
\deleted{\lambda\in H_x,\quad u\in M\Longrightarrow\lambda u=\lambda}
\]
application
\[
u\longmapsto fu:\ H_e\longrightarrow H_x
\]
définie par un $f\in F(x):e\longrightarrow x$ ne dépend pas du choix de $f$, donc
\[
fu=gu\quad\text{pour }\begin{aligned}&u\in H_e,\\& f,g:e\rightrightarrows x.\end{aligned}
\]

% Printed form upside down relative to the mathematics.
% letterhead
{\footnotesize
\begin{tabular}{@{}l@{\quad}l@{}}
\deleted{MINISTERE DE L'EDUCATION NATIONALE} & \deleted{AVIGNON, le}\\
\deleted{ACADEMIE D'AIX-MARSEILLE} &\\
\hfil\rule{1.5cm}{0.4pt}\hfil &\\
\deleted{Inspection Académique} & \deleted{L'Inspecteur d'académie,}\\
\deleted{de Vaucluse} & \deleted{directeur des services départementaux}\\
\hfil\rule{1.5cm}{0.4pt}\hfil & \deleted{de l'Education Nationale de Vaucluse,}\\
\deleted{Tél. : 90 82 99 76} &\\
\deleted{Référence : Secrétariat n\textsuperscript{o}} & \deleted{à M}
\end{tabular}}
% end letterhead
% archive sheet number: 51
